\section{Scheduling Module}

\subsection{Module}
The Scheduling Module manages recurring and one-time ride schedules for drivers and riders.
It allows users to automate trip creation, define weekly patterns, and ensures that recurring schedules generate valid trip entries.  
This module integrates directly with the Trip Module and the Notification Module.

\paragraph{Frontend}
Implemented in the Expo app, the frontend provides UI controls for defining recurring or one-time schedules.

Key capabilities:
\begin{itemize}
    \item Time picker for selecting departure time
    \item Day-of-week toggles for recurring schedules
    \item Interface to view and cancel existing schedules
    \item Integration with tRPC hooks such as \texttt{useCreateSchedule}, \texttt{useGetUserSchedules}
    \item Form validation for time ranges and required fields
\end{itemize}

\paragraph{Backend}
The backend implements scheduling logic via tRPC.

Responsibilities include:
\begin{itemize}
    \item Parsing user-defined recurrence patterns
    \item Generating future trips based on schedules
    \item Ensuring schedules do not conflict with existing trips
    \item Validating schedule timing (future date, valid days)
    \item Exposing procedures such as \texttt{createSchedule}, \texttt{cancelSchedule}
\end{itemize}

\paragraph{Data}
The module persists schedule definitions and generated trip associations.

Tables:
\begin{itemize}
    \item \texttt{schedules} --- recurring schedule definitions (days, time, userId)
    \item \texttt{schedule\_instances} --- maps schedules to generated trips
    \item \texttt{trips} --- created by this module but stored by Trip Module
\end{itemize}

\subsection{Uses}
\begin{itemize}
    \item Trip Module --- creates trip entries from schedules
    \item User Profile Module --- determines valid schedule types
    \item Notification Module --- sends reminders for upcoming rides
\end{itemize}

\subsection{Syntax}

\paragraph{Exported Constants}
\begin{itemize}
    \item \texttt{MAX\_RECURRING\_YEARS} = 1 --- Maximum duration in years to generate future instances.
    \item \texttt{VALID\_DAYS} = \{Monday, Tuesday, Wednesday, Thursday, Friday, Saturday, Sunday\}
\end{itemize}

\paragraph{Exported Access Programs}
\par\vspace{0.5em}

\begin{center}
\resizebox{\textwidth}{!}{
\begin{tabular}{l l l l}
\toprule
Name & In & Out & Exceptions \\
\midrule
\texttt{createSchedule} & userId: $\mathbb{N}$, data: ScheduleData & ScheduleRecord & ValidationError \\
\texttt{getUserSchedules} & userId: $\mathbb{N}$ & ScheduleList & NotFoundError \\
\texttt{cancelSchedule} & userId: $\mathbb{N}$, scheduleId: $\mathbb{N}$ & Confirmation & NotFoundError, UnauthorizedError \\
\texttt{generateFutureTrips} & scheduleId: $\mathbb{N}$ & TripList & AlgorithmError \\
\bottomrule
\end{tabular}
}
\end{center}


\subsection{Semantics}

\paragraph{State Variables}
\begin{itemize}
    \item \texttt{Schedules}: Map of ($\mathbb{N} \rightarrow$ ScheduleRecord) --- The collection of all schedule definitions indexed by \texttt{scheduleId}.
    \item \texttt{ActiveRules}: Map of ($\mathbb{N} \rightarrow$ string) --- The collection of parsed recurrence rules (e.g., cron strings) indexed by \texttt{scheduleId}.
\end{itemize}

\paragraph{Environment Variables}
\begin{itemize}
    \item System clock and device timezone for temporal calculation and validation.
    \item Persistence layer connection (Database).
\end{itemize}

\paragraph{Assumptions}
\begin{itemize}
    \item The user is authenticated, and their \texttt{userId} is verified by the Authentication Module prior to invocation.
    \item The database schema is fully migrated and accessible.
\end{itemize}

\paragraph{Access Routine Semantics}

\texttt{createSchedule}(userId, data):
\begin{itemize}
    \item \textbf{transition:} Generates a unique \texttt{id}. Formats \texttt{newSchedule} combining \texttt{id}, \texttt{userId}, and \texttt{data}. Updates state: 
    \begin{itemize}
        \item $\texttt{Schedules[id]} := \texttt{newSchedule}$
        \item $\texttt{ActiveRules[id]} := \texttt{parseRecurrenceRule(data)}$
    \end{itemize}
    \item \textbf{output:} Returns \texttt{Schedules[id]}.
    \item \textbf{exception:} \texttt{ValidationError} if \texttt{data.isRecurring = true} $\land$ \texttt{data.daysOfWeek} $= \emptyset$, or if \texttt{data.time} is in the past relative to the system clock.
\end{itemize}

\texttt{getUserSchedules}(userId):
\begin{itemize}
    \item \textbf{transition:} None.
    \item \textbf{output:} Returns the sequence $\{s \in \texttt{Schedules} \mid s.\text{userId} = \text{userId}\}$.
    \item \textbf{exception:} \texttt{NotFoundError} if the resulting sequence is empty.
\end{itemize}

\texttt{cancelSchedule}(userId, scheduleId):
\begin{itemize}
    \item \textbf{transition:} Removes \texttt{scheduleId} from \texttt{Schedules} and \texttt{ActiveRules}. Notifies dependent database constraints to cascade status changes.
    \item \textbf{output:} Returns \texttt{true} upon successful cancellation.
    \item \textbf{exception:} \texttt{NotFoundError} if $\texttt{scheduleId} \notin \texttt{Schedules}$. \texttt{UnauthorizedError} if $\texttt{Schedules[scheduleId].userId} \neq \text{userId}$.
\end{itemize}

\texttt{generateFutureTrips}(scheduleId):
\begin{itemize}
    \item \textbf{transition:} 
    \begin{itemize}
        \item Retrieves rule $r = \texttt{ActiveRules[scheduleId]}$.
        \item Uses \texttt{nextOccurrence} to compute dates.
        \item Generates \texttt{TripRecord} entries via the Trip Module up to \texttt{MAX\_RECURRING\_YEARS}.
    \end{itemize}
    \item \textbf{output:} Returns the generated sequence of \texttt{TripRecord} entities.
    \item \textbf{exception:} \texttt{AlgorithmError} if $\texttt{scheduleId} \notin \texttt{ActiveRules}$ or if the string rule fails to parse into valid future dates.
\end{itemize}

\paragraph{Local Functions}
\begin{itemize}
    \item \texttt{parseRecurrenceRule(data: ScheduleData): string} --- Converts schedule days and time into a cron-formatted string or scheduling equivalent.
    \item \texttt{nextOccurrence(rule: string, after: $\mathbb{R}$): $\mathbb{R}$} --- Computes the next valid Unix timestamp given a recurrence rule and a baseline timestamp.
\end{itemize}